% context: markscheme (official solutions) for 1-2CW_Velocity_v2
% differs from 1-2CW_Velocity_MS only in Q6: v2 drops the metric second
%   parts of 6(b) and 6(d), kept here as a teacher aside
% topic: average speed, unit conversion between ft/s, mi/h, m/s, km/h
% convention: "=" joins exact equalities; "\approx" introduces a rounded value
\documentclass[12pt, twoside]{article}
\input{../preamble}
\usepackage{float}
\setlength{\parskip}{0.3\baselineskip}

\title{Physics}
\author{Chris Huson}
\date{September 2026}

\fancyhead[LE]{\thepage}
\fancyhead[RO]{\thepage}
\fancyhead[LO]{La Scuola d'Italia / Huson / Physics \\* 17 September 2026}

\begin{document}

\subsubsection*{1.2 Classwork Solutions: Speed, Velocity, and Unit Conversion}

\emph{Notation: \,$=$\, joins values that are exactly equal; \,$\approx$\, introduces a
rounded value. Students should write it the same way --- once you round the calculator
display, the equals sign is no longer true. Each conversion runs as one chain of unit
factors, so there is nothing to round in the middle.}

\begin{enumerate}[itemsep=0.2cm]

\item A student walks 24.0 m in 18.5 s.
  \begin{enumerate}[itemsep=0.2cm]
    \item $v = \dfrac{d}{t} = \dfrac{24.0~\text{m}}{18.5~\text{s}} = 1.2972\ldots \approx \mathbf{1.30~\textbf{m/s}}$
    \item \textbf{Three significant figures.} Both measurements have three, and a quotient
      carries no more than the least precise measurement that went into it. Writing 1.29729
      would claim a precision the stopwatch and the tape measure do not have.
  \end{enumerate}

\item Convert each speed.
  \begin{enumerate}[itemsep=0.2cm]
    \item $30.0~\dfrac{\text{km}}{\text{h}} \times \dfrac{1000~\text{m}}{1~\text{km}}
      \times \dfrac{1~\text{h}}{3600~\text{s}}
      = \dfrac{30{,}000}{3600}~\dfrac{\text{m}}{\text{s}} = 8.333\ldots
      \approx \mathbf{8.33~\textbf{m/s}}$

      \emph{The km cancel and the h cancel, leaving m/s. A useful benchmark: dividing
      km/h by 3.6 gives m/s.}
    \item $25.0~\dfrac{\text{m}}{\text{s}} \times \dfrac{1~\text{km}}{1000~\text{m}}
      \times \dfrac{3600~\text{s}}{1~\text{h}} = \mathbf{90.0~\textbf{km/h}}$
      \quad (exact --- nothing to round)
  \end{enumerate}

\item A fastball at 95 mi/h.
  \begin{enumerate}[itemsep=0.2cm]
    \item $95~\dfrac{\text{mi}}{\text{h}} \times \dfrac{5280~\text{ft}}{1~\text{mi}}
      \times \dfrac{1~\text{h}}{3600~\text{s}}
      = \dfrac{501{,}600}{3600}~\dfrac{\text{ft}}{\text{s}} = 139.33\ldots
      \approx \mathbf{139~\textbf{ft/s}}$
    \item Chain all the way from the original figure in one expression:

      $95~\dfrac{\text{mi}}{\text{h}} \times \dfrac{5280~\text{ft}}{1~\text{mi}}
      \times \dfrac{1~\text{m}}{3.28~\text{ft}} \times \dfrac{1~\text{h}}{3600~\text{s}}
      = \dfrac{95 \times 5280}{3.28 \times 3600}~\dfrac{\text{m}}{\text{s}}
      = 42.48\ldots \approx \mathbf{42.5~\textbf{m/s}}$

      \emph{Write the whole conversion as one product and evaluate once. If you continue
      from part (a) instead, carry the full 139.33 ft/s --- restarting from the rounded 139
      gives 42.4, wrong in the last figure. Rounding mid-chain is the commonest way to lose
      an accuracy mark.}
  \end{enumerate}

\item Put all three in m/s before comparing.
  \begin{itemize}[itemsep=0.25cm]
    \item Car: $60~\dfrac{\text{mi}}{\text{h}} \times \dfrac{5280~\text{ft}}{1~\text{mi}}
      \times \dfrac{1~\text{m}}{3.28~\text{ft}} \times \dfrac{1~\text{h}}{3600~\text{s}}
      = \dfrac{60 \times 5280}{3.28 \times 3600}~\dfrac{\text{m}}{\text{s}}
      = 26.829\ldots \approx \mathbf{26.8~\textbf{m/s}}$
    \item Train: $95~\dfrac{\text{km}}{\text{h}} \times \dfrac{1000~\text{m}}{1~\text{km}}
      \times \dfrac{1~\text{h}}{3600~\text{s}} = 26.39\ldots \approx \mathbf{26.4~\textbf{m/s}}$
    \item Motorcycle: $\mathbf{27~\textbf{m/s}}$ (already in m/s)
  \end{itemize}
  \vspace{0.15cm}
  At three figures the order reads train $<$ car $<$ motorcycle --- \textbf{but that
  ordering is not supported by the data.}

  \emph{Every speed here is given to two significant figures at best --- and ``60 mi/h''
  may be only one, since the trailing zero is ambiguous. At two figures the three speeds
  are 27, 26 and 27 m/s, so the defensible answer is: \textbf{the train is the slowest,
  and the car and the motorcycle cannot be told apart.} The 0.2 m/s between 26.8 and 27 is
  smaller than the uncertainty already in the data; using the third figure to break that
  tie claims precision the problem never gave. Give full credit for a correctly converted
  list, and full credit again --- the real lesson --- for saying the ranking cannot be
  completed.}

\item Sound travels 343 m/s; the delay is 4.0 s.

  $d = vt = (343~\text{m/s})(4.0~\text{s}) = 1372~\text{m} \approx \mathbf{1.4 \times 10^{3}~\textbf{m}}$
  \quad (2 s.f., because 4.0 s has two)

  $1372~\text{m} \times \dfrac{3.28~\text{ft}}{1~\text{m}} \times \dfrac{1~\text{mi}}{5280~\text{ft}}
  = 0.8523\ldots \approx \mathbf{0.85~\textbf{mi}}$

  \emph{This is the ``count the seconds'' rule of thumb: roughly 1 km every 3 seconds,
  or about 1 mile every 5 seconds.}

\item Climbing 5 floors, 18 steps each, 8.0 in per step.
  \begin{enumerate}[itemsep=0.2cm]
    \item $5~\text{floors} \times \dfrac{18~\text{steps}}{1~\text{floor}} = \mathbf{90~\textbf{steps}}$
    \item $90~\text{steps} \times \dfrac{8.0~\text{in}}{1~\text{step}}
      \times \dfrac{1~\text{ft}}{12~\text{in}} = \mathbf{60~\textbf{ft}}$

      \emph{The step counts are exact, so the 8.0 in riser sets the precision: two
      significant figures. In metric, if it comes up:
      $60~\text{ft} \times \frac{1~\text{m}}{3.28~\text{ft}} \approx 18$ m.}
    \item $t = \dfrac{90~\text{steps}}{1.2~\text{steps/s}} = \mathbf{75~\textbf{s}} = \mathbf{1~\textbf{min}~15~\textbf{s}}$
    \item $v = \dfrac{60~\text{ft}}{75~\text{s}} \times \dfrac{60~\text{s}}{1~\text{min}}
      = \mathbf{48~\textbf{ft/min}}$

      \emph{One chain from the 60 ft in 75 s --- there is no need to work out ft/s first
      and convert again. In metric, if it comes up:
      $\frac{60~\text{ft}}{75~\text{s}} \times \frac{1~\text{m}}{3.28~\text{ft}} \approx 0.24$ m/s.
      Note this is the \emph{vertical} speed; the distance walked along the stairs is longer,
      since each step also moves you forward.}
  \end{enumerate}

\item \textbf{Challenge.} The conversion itself is right:

  $30~\dfrac{\text{mi}}{\text{h}} \times \dfrac{1.61~\text{km}}{1~\text{mi}} = 48.3~\text{km/h}
  \approx \mathbf{48~\textbf{km/h}}$

  What is wrong is the number of digits. The speedometer reads 30 mi/h --- two significant
  figures --- so the converted speed can carry at most two. The seven digits in 48.28032
  came out of the calculator, not the measurement; they claim the car's speed is known to a
  hundred-thousandth of a km/h. (The conversion factor is rounded too: 1 mi = 1.609344 km.)

  \emph{Converting a measurement never makes it more precise --- it only restates the
  same precision in a different unit.}

\end{enumerate}
\end{document}
